\chapter{مداخل الآلة: كيف توصَل بها العيون والآذان والمستشعرات الأخرى}
\chaptermark{مداخل الآلة}
\label{sec:sensors}
\begin{chaptergoals}
\item كيف تعمل كل حاسة محوِّلًا من مستقبِل بوابة وضبط للكسب ومرمّز للنبضات؛
\item لماذا يحدّ ضجيج الطلقات للفوتونات الرؤيةَ في الظلام، ولماذا نرى في الغسق أبطأ؛
\item كيف يحشر الضبط التلقائي للكسب مدى هائلًا في النبضات، فتنقل المستشعرات التغيرات؛
\item كيف تضغط العين صورتها نحو مئة مرة، وكيف تعمل الأذن بنكَ مرشحات.
\end{chaptergoals}

\section{المستشعر محوِّلًا}

في الشكل~\ref{fig:machine} كانت البيانات تدخل الآلة من اليسار، من كتلة ”المستشعرات“. فلنفتح الآن هذه الكتلة. كل عضو حسّ عند المهندس \emph{محوِّل} له طور دخل تماثلي وفي آخره محوِّل تماثلي رقمي، إلا أن الأرقام في المخرج ليست أعدادًا، بل نبضات.

المخطط واحد في كل أعضاء الحس (الشكل~\ref{fig:transducer}). كمية فيزيائية، ضوءًا كانت أو ضغطًا أو اهتزازًا أو جزيئًا، تؤثر في بروتين مستقبِل في غشاء الخلية الحسية. يعمل هذا البروتين عمل البوابة: يفتح قناة أيونية بنفسه أو عبر سلسلة قصيرة من التفاعلات. فينشأ \emph{تيار المستقبِل}، ومعه جهد تماثلي على الغشاء. ثم يمر الجهد بضبط الكسب ويصل إلى مرمّز النبضات، وهو العصبون المكامل المعروف لنا~\eqref{eq:glutneuron}، الذي يحوّل المستوى إلى تردد وإلى لحظات نبضات. والمخرج واحد في الغالب: العصبونات الحسية في العين والأذن والشم والجلد تطلق الغلوتامات، فتتصل المداخل مباشرة بناقل \nt{GLUT}.

\begin{figure}[t]
\centering
\begin{tikzpicture}[>=Stealth,
  blk/.style={draw,very thick,rounded corners=2pt,align=center,font=\footnotesize,minimum height=1.0cm,text width=1.9cm,fill=blue!5}]
\node[align=center,font=\small] (P) at (0.1,0) {ضوء، صوت،\\ضغط،\\جزيء};
\node[blk] (R) at (3.0,0) {المستقبِل\\البوابة};
\node[blk] (G) at (6.5,0) {ضبط\\الكسب};
\node[blk] (E) at (10.0,0) {مرمّز\\النبضات};
\node[align=center,font=\small] (B) at (13.4,0) {ناقل\\\nt{GLUT}};
\draw[->,very thick] (P) -- (R);
\foreach \ic/\xx in {sun/-1.1,speaker/-0.3,press/0.5,molecule/1.3}{\pic[scale=0.85] at (\xx,1.05) {\ic};}
\draw[->,very thick] (R) -- node[above,font=\scriptsize] {تيار} (G);
\draw[->,very thick] (G) -- node[above,font=\scriptsize] {مستوى} (E);
\draw[->,very thick,red!70!black] (E) -- node[above,font=\scriptsize,black] {نبضات} (B);
\draw[decorate,decoration={brace,amplitude=5pt,mirror},gray!60!black] (1.8,-0.8) -- (7.7,-0.8) node[midway,below=5pt,font=\scriptsize,black] {الجزء التماثلي};
\draw[decorate,decoration={brace,amplitude=5pt,mirror},gray!60!black] (8.8,-0.8) -- (11.2,-0.8) node[midway,below=5pt,font=\scriptsize,black] {نبضات};
\end{tikzpicture}
\caption{أي عضو حسّ سلسلةً من التحويلات. يفتح البروتين المستقبِل قناة فيعطي تيارًا؛ ويكيّفه ضبط الكسب مع الظروف؛ ويحوّل المرمّز~\eqref{eq:glutneuron} المستوى إلى نبضات تخرج إلى ناقل \nt{GLUT}. في العين والأذن لا تُصدر الأطوار الأولى نبضات أصلًا: تبقى الإشارة تماثلية إلى أن يلزم نقلها بعيدًا.}
\label{fig:transducer}
\end{figure}

تفصيل يعرفه مهندس الإلكترونيات جيدًا. في العين والأذن لا تُصدر الخلايا الأولى نبضات إطلاقًا: إنها تنقل إلى جاراتها جهدًا متدرجًا، كطور تماثلي على لوحة دارة. ولا تظهر النبضات إلا حيث يجب نقل الإشارة بعيدًا. فالإشارة التماثلية تخمد وتتلوث بالضجيج على مسافة ملليمترات، أما النبضة، كما رأينا في الفصل~\ref{sec:neurontypes}، فتصل إلى نهاية السلك بالارتفاع نفسه. والرقمنة تجري حيث يبدأ الخط الطويل.

\section{عند حدود الفيزياء}

تعمل مستشعرات الدماغ عند الحدود الفيزيائية للحساسية. مستشعر الضوء في الظلام يستجيب \emph{لفوتون واحد}: امتصاص جسيم ضوء واحد يعطي تيارًا نحو بيكوأمبير، يمكن تمييزه من ضجيج المستشعر نفسه~\cite{Baylor1979}. ومستشعر الصوت يلاحظ إزاحة حزمته الحساسة بأقل من نانومتر~\cite{Hudspeth1989}.

حين يقل الضوء لا يحدّ الدقةَ المستشعرُ، بل الضوء نفسه: تصل الفوتونات عشوائيًا في تدفق بواسوني. إذا دخل المستشعر خلال زمن التجميع $T$ في المتوسط $N=\Phi T$ فوتونًا، تذبذب عددها بمقدار $\sqrt N$. ولكي تُلاحظ زيادة $\Delta N$ يجب أن تكون أكبر من هذا التذبذب بعدة مرات، فتكون أصغر نسبة لتغير السطوع يمكن تمييزها
\begin{empheq}[box=\keybox]{equation}
\frac{\Delta I}{I} \;\approx\; \frac{k}{\sqrt{\Phi\,T}} .
\label{eq:photonnoise}
\end{empheq}
هذه هي الصيغة نفسها التي رأيناها لعدد النبضات~\eqref{eq:rateerror}: ضجيج الطلقات للفوتونات وضجيج الطلقات للنبضات يخضعان لقانون واحد. وفي الظلام لا تملك العين إلا طريقة واحدة لتحسين الدقة، هي أن تجمع الفوتونات مدة أطول، وزمن التجميع يطول فعلًا حتى أعشار الثانية. لهذا نرى في الغسق أبطأ.

\section{المدى الديناميكي: ضبط الكسب}

أصعب مسألة هندسية في المستشعرات هي المدى. تتغير الإضاءة من ليلة مرصعة بالنجوم إلى ثلج تحت الشمس بنحو عشر رتب عشرية، وشدة الصوت من عتبة السمع إلى عتبة الألم باثنتي عشرة رتبة. أما تردد النبضات فيتغير من الصفر إلى بضع مئات في الثانية، أي بأقل من ثلاث رتب. ولا يمكن حشر دخل كهذا في مخرج كهذا خطيًا.

والحل هو حل أي مستقبِل لاسلكي: \emph{الضبط التلقائي للكسب}. توصف استجابات الخلايا الحسية في العين وصفًا جيدًا بالصيغة
\begin{empheq}[box=\keybox]{equation}
r \;=\; r_{\max}\,\frac{I}{I+\sigma} ,
\label{eq:nakarushton}
\end{empheq}
حيث $\sigma$ السطوع الذي تبلغ عنده الاستجابة نصف قيمتها العظمى~\cite{NakaRushton1966}. والصيغة~\eqref{eq:nakarushton} وحدها لا تغطي إلا رتبتين أو ثلاثًا. لكن $\sigma$ ليست ثابتة: عند العمل طويلًا على خلفية ساطعة تزداد تبعًا لمتوسط السطوع. فينزاح منحنى الاستجابة على المحور اللوغاريتمي للسطوع، ويضع جزأه الشديد الانحدار في كل مرة حيث يقع الدخل الآن (الشكل~\ref{fig:adaptation}). هذه هي القسمة نفسها على النشاط الكلي التي رأيناها في التثبيط التحويلي في الفصل~\ref{sec:gaba}~\cite{Carandini2012}، إلا أن المقام هنا متوسط السطوع خلال الثواني الأخيرة.

\begin{figure}[t]
\centering
\begin{tikzpicture}[>=Stealth,x=1.25cm,y=2.8cm]
\draw[->,gray!70!black] (-2,0) -- (6.4,0) node[below left,font=\small,black] {$\log_{10} I$};
\draw[->,gray!70!black] (-2,0) -- (-2,1.15) node[left,font=\small,black] {$r/r_{\max}$};
\foreach \u in {-2,0,2,4,6}{\draw[gray!60!black] (\u,-0.02) -- (\u,0.02); \node[below,font=\scriptsize] at (\u,0) {$\u$};}
\draw[densely dashed,gray!60!black] (-2,0.5) -- (6.2,0.5);
\foreach \s/\c/\lab in {0/blue!60!black/{خلفية مظلمة},2/green!45!black/{متوسطة},4/red!70!black/{ساطعة}}{
  \pgfmathsetmacro{\lo}{max(-2,\s-3)}
  \draw[very thick,\c] (-2,0) -- (\lo,0) plot[domain=\lo:6.2,samples=80] (\x,{1/(1+10^(\s-\x))});
  \node[font=\scriptsize,\c,anchor=west] at (\s+0.15,0.42) {\lab};}
\foreach \s/\ic in {0/moon,2/cloud,4/sun}{\pic at (\s+0.6,0.64) {\ic};}
\end{tikzpicture}
\caption{ضبط كسب مستشعر الضوء بحسب~\eqref{eq:nakarushton}. كل منحنى يغطي رتبتين أو ثلاثًا من السطوع؛ وعند الانتقال إلى خلفية أسطع تزداد $\sigma$ فينزاح المنحنى على المحور اللوغاريتمي. والمنحنيات المنزاحة معًا تغطي المدى كله.}
\label{fig:adaptation}
\end{figure}

لهذا الضبط ثمن نعرفه من الفصل~\ref{sec:code}: لا يُبلغ المستشعر عن السطوع، بل عن السطوع \emph{بالنسبة إلى الخلفية}. الدخل الثابت يكفّ تدريجيًا عن استدعاء الاستجابة، أما كل تغيير فيستدعيها. فمداخل الآلة تتحدث منذ البداية عن التغيرات لا عن المستويات، وفي هذا الاقتصاد نفسه في النبضات الذي في القضية~\ref{prop:economy}~\cite{Barlow1961}. ومن هنا أيضًا مستشعرات التشغيل والإطفاء في الفصل~\ref{sec:glutcomputer}~\cite{Schiller1992}: بعضها يُبلغ أن الضوء ازداد، وبعضها أنه خفت.

وقد كرر المهندسون هذه الحيلة في \emph{كاميرات الأحداث}. لا تلتقط هذه الكاميرا إطارات وفق ساعة: كل بكسل فيها يُصدر بنفسه، بلا ساعة مشتركة، حدثًا ”أسطع“ أو ”أعتم“ حين يتغير لوغاريتم السطوع بنسبة محددة~\cite{Lichtsteiner2008}. هذه بالضبط شفرة السلكين للتشغيل والإطفاء من الفصل~\ref{sec:glutcomputer} مجسَّدة في السيليكون، وهي تمنح الكاميرا مدى وسرعة لا تبلغهما المصفوفة العادية~\cite{Gallego2022}.

\section{العين: ضغط حتى يتسع له السلك}

في العين نحو $10^8$ مستشعر ضوء~\cite{Curcio1990}، أما عصبونات المخرج التي تنقل الصورة إلى الدماغ فنحو مليون. وقبل أن تغادر الإشارة العين تُضغط نحو مئة مرة. وأهم حيل الضغط معروفة لنا. كل عصبون مخرج يستجيب للفرق بين السطوع في بقعة صغيرة والسطوع حولها، وهذا طرح الجيران عبر التثبيط كما في الفصل~\ref{sec:gaba}. المساحات المستوية من الصورة لا تكاد تكلف شيئًا، أما الحواف والتغيرات فتُنقل. وعصبونات المخرج المختلفة متخصصة: بعضها يُبلغ عن الإضاءة، وبعضها عن الإعتام، وبعضها عن الحركة في اتجاه معيّن (الشكل~\ref{fig:direction})؛ ولدى الفأر أُحصي أكثر من ثلاثين قناة مخرج كهذه~\cite{Baden2016}.

فما عرض نطاق الكابل الناتج؟ لدى خنزير غينيا قيس من تسجيلات عصبونات المخرج نحو $875$ ألف بت في الثانية لنحو $\sim10^5$ من هذه العصبونات؛ والتحويل إلى مليون عصبون مخرج لدى الإنسان يعطي رتبة $10^7$ بت في الثانية~\cite{Koch2006}، أي ما يعادل كابل شبكة قديمًا بسرعة عشرة ميغابت. وبتردد متوسط يبلغ بضع نبضات في الثانية لكل عصبون يكون ذلك بضع بتات لكل نبضة، كما خرج لنا في الفصل~\ref{sec:code}.

\section{الأذن: بنك مرشحات}

تحل الأذن مسألة أخرى: يجب تحليل الصوت إلى تردداته. كان المهندس سيضع بنكًا من مرشحات تمرير النطاق، والأذن تفعل ذلك بالضبط ميكانيكيًا. يجري الاهتزاز الصوتي على حاجز مرن تتغير خصائصه تدريجيًا على طوله، فيرنّ كل جزء منه على تردده. والمستشعرات الجالسة على جزء ما لا تستجيب إلا لنطاقها (الشكل~\ref{fig:filterbank}). وتتوزع الترددات على المواضع توزيعًا لوغاريتميًا تقريبًا: كل أوكتاف ينال حصة متقاربة من المستشعرات. ورقم المستشعر الذي أطلق هو التردد، وهذه شفرة الموضع من الفصل~\ref{sec:code}.

\begin{figure}[t]
\centering
\begin{tikzpicture}[>=Stealth,x=1.35cm,y=2.2cm]
\draw[->,gray!70!black] (-0.3,0) -- (7.6,0) node[above left,font=\small,black] {التردد، \foreignlanguage{english}{Hz}};
\draw[->,gray!70!black] (-0.3,0) -- (-0.3,1.15) node[left,font=\small,black] {الاستجابة};
\foreach \k/\f in {0/125,1/250,2/500,3/1000,4/2000,5/4000,6/8000}{
  \draw[gray!60!black] (\k,-0.02) -- (\k,0.02); \node[below,font=\scriptsize] at (\k,0) {\f};
  \draw[thick,blue!60!black] plot[domain={max(\k-1.2,-0.3)}:{\k+0.7},samples=60] (\x,{1/(1+((\x-\k)/(\x<\k?0.45:0.2))^4)});}
% a piano keyboard: seven white keys per octave, like the octaves on the axis
\foreach \i in {0,...,48}{\draw[gray!50!black,fill=white,line width=0.3pt] ({\i/7},-0.44) rectangle ({(\i+1)/7},-0.26);}
\foreach \o in {0,...,6}{\foreach \j in {0,1,3,4,5}{\fill[black] ({\o+(\j+1)/7-0.035},-0.35) rectangle ({\o+(\j+1)/7+0.035},-0.26);}}
\end{tikzpicture}
\caption{الأذن بنكًا من مرشحات تمرير النطاق، مخطط. كل منحنى استجابةُ مستشعرات جزء واحد لصوت بترددات مختلفة؛ والترددات المركزية متباعدة بأوكتاف، كمفاتيح على مقياس لوغاريتمي. للمرشحات الحقيقية منحدر شديد في جهة الترددات العالية ومنحدر لطيف في جهة المنخفضة. وفي الأسفل لوحة مفاتيح: عليها، كما في الأذن، ينال كل أوكتاف المساحة نفسها.}
\label{fig:filterbank}
\end{figure}

مرنانات الأذن ليست منفعلة. فبعض المستشعرات يهزّ الحاجز بنفسه بإيقاع الصوت، ويضخّم الاهتزازات الضعيفة آلاف المرات في السعة (بمقدار $66\text{--}76\,\text{dB}$)، ومع ازدياد الشدة ينخفض الكسب~\cite{Ruggero1997,Robles2001}. هذا عند المهندس مضخّم بتغذية راجعة موجبة موضوع عند حافة الإثارة الذاتية تمامًا، وهي الحياة نفسها على الحافة التي رأيناها في شبكة الفصل~\ref{sec:gaba}: عند الحافة يكون النظام حساسًا للإشارات الصغيرة على نحو خاص. وضغط الشدة يؤديه المضخّم نفسه: الهادئ يُضخَّم كثيرًا، والصاخب قليلًا.

وللأذن شفرة ثانية. في الترددات المنخفضة تسير نبضات العصبونات في طور مع الصوت: يطلق العصبون عند جزء معيّن من كل موجة، وإن لم يطلق عند كل موجة. لدى خنزير غينيا يبدأ التقيد بالطور بالضعف قرب $600\,\text{Hz}$، ويبقى ملحوظًا حتى نحو $3.5\,\text{kHz}$~\cite{PalmerRussell1986}. هكذا يُحفظ في لحظات النبضات التوقيت الدقيق للصوت، ومنه الفرق بين وصول الصوت إلى الأذنين، الذي تجد به شبكة الفصل~\ref{sec:glutcomputer} الاتجاه~\cite{Grothe2010}. وفي الترددات العالية لا تلحق النبضات بالموجة، فلا تبقى إلا شفرة الموضع. فالأذن الواحدة تستعمل شفرتي الفصل~\ref{sec:code} كلتيهما: التوقيت حيث تلحق العصبونات، والموضع حيث لا تلحق.

\section{بقية المستشعرات}

\begin{table}[t]
\centering
\footnotesize
\setlength{\tabcolsep}{3pt}
\renewcommand{\arraystretch}{1.15}
\begin{tabular}{>{\raggedright}p{1.9cm}>{\raggedright}p{2.6cm}>{\raggedright}p{4.0cm}>{\raggedright\arraybackslash}p{4.6cm}}
\toprule
الحاسة & ماذا تقيس & كيف بُني المدخل & الحيلة من الكتاب \\
\midrule
البصر & الضوء & $\sim10^8$ مستشعر، ضغط $\sim100$ مرة، $\sim10^7$ بت في الثانية & ضبط الكسب، طرح الجيران، سلكان، اتجاه الحركة \\
السمع & ضغط الهواء & بنك مرشحات ميكانيكي، مضخّم فعّال & شفرة الموضع وشفرة التوقيت، التأخيرات \\
اللمس & الضغط، الاهتزاز & مستشعرات سريعة التعوّد وبطيئته & زوج ”مرشح تمرير عالٍ ومرشح تمرير منخفض“ \\
الحرارة & الدفء & مستشعرات منفصلة للدفء وللبرد & شفرة السلكين \\
الشم & الجزيئات & $\sim400$ نوع من المستقبِلات، ولكل مستشعر نوع واحد & شفرة تركيبية: الرائحة مجموعة الأنواع التي أطلقت \\
الذوق & مواد الطعام & خمس صفات: حلو، مرّ، مالح، حامض، أومامي & خطوط موسومة؛ بعض الخلايا يطلق \foreignlanguage{english}{ATP} أو السيروتونين لا الغلوتامات \\
وضع الجسم & تمدد العضلات & مستشعرات الطول وسرعة التمدد & تغذية راجعة لمنظّم الحركة \\
\bottomrule
\end{tabular}
\caption{كيف توصَل المستشعرات. كل الأجهزة في العمود الثالث مقيسة؛ والعمود الأخير يربطها بحيل الفصول السابقة.}
\label{tab:sensors}
\end{table}

بقية الحواس مجموعة في الجدول~\ref{tab:sensors}، وفي كل سطر منه تقريبًا نتعرف على حيلة من الفصول السابقة. يستعمل اللمس زوجًا من المرشحات: بعض مستشعرات الضغط يستجيب للمسة ثم يصمت فورًا، وبعضها يُبقي استجابته ما دام الضغط قائمًا. وتنقل الحرارةَ سلكان، أحدهما للدفء والآخر للبرد. وأغرب المداخل الشم. لدى الإنسان نحو أربعمئة نوع من المستقبِلات الشمية، ولا يحمل كل عصبون شمي إلا نوعًا واحدًا~\cite{Mombaerts2004}. والرائحة لا تنشّط نوعًا واحدًا بل مجموعة، والرسالة هي \emph{أيّ} الأنواع أطلقت. هذا عند المهندس متجه ثنائي طوله نحو أربعمئة، وعدد قيمه الممكنة فلكي.

\begin{keyblock}
\begin{proposition}[مداخل الآلة]
\label{prop:sensors}
كل حاسة محوِّل يعمل عند الحد الفيزيائي للحساسية، ويضغط مدى ديناميكيًا هائلًا بالضبط التلقائي للكسب، وينقل على ناقل \nt{GLUT} \emph{التغيرات} قبل كل شيء. ولهذا تستعمل المستشعرات حيل الآلة نفسها: شفرة السلكين، وشفرة الموضع، وشفرة التوقيت، والقسمة على الخلفية. بنية المستشعرات مقيسة؛ أما أن شفراتها مضبوطة على أكبر قدر من المعلومات لكل نبضة ففرضية راسخة الأسس.
\end{proposition}
\end{keyblock}

تعمل المداخل، ككل شيء آخر، بلا تزامن. كل مستشعر يُصدر نبضة حين يكون لديه ما يبلّغ عنه، والحواس المختلفة تصل بتأخيرات مختلفة: الصوت بعد ملّي ثوانٍ، والضوء بعد عشرات الملّي ثانية. أما كيف تجمعها الآلة في صورة واحدة للعالم من غير ساعة مشتركة، فهي إحدى مسائل التنسيق في الزمن من الفصل~\ref{sec:synthesis}.

\section{تمارين}

\begin{exercise}[\lvlA]
\label{ex:11:1}
\emph{كم نجمع من الفوتونات.} في الغسق يدخل المستشعرَ $100$ فوتون في الثانية؛ والزيادة ملحوظة إذا كانت أكبر من التذبذب~\eqref{eq:photonnoise} بـ$3$ مرات. (أ)~كم يحتاج مستشعر واحد من الوقت ليلاحظ تغير السطوع بـ$10\,\%$؟ (ب)~كم مستشعرًا يجب جمعها لينجز ذلك في $0.2\,\text{s}$، وبكم مرة تنخفض الدقة المكانية على المحور؟
\end{exercise}

\begin{exercise}[\lvlA]
\label{ex:11:2}
\emph{كم بتًا في النبضة.} (أ)~تنقل العين $\sim10^7$ بت في الثانية عبر $10^6$ عصبون مخرج بتردد من $3$ إلى $5$ نبضات في الثانية. ما النسبة التي تستعملها من السعة~\eqref{eq:spikeentropy} عند $\Delta t=1\ms$؟ (ب)~رائحة تشغّل $20$ من $\approx400$ نوع مستقبِل. كم بتًا تحمل هذه المجموعة، وكم نوعًا مشتركًا في المتوسط بين رائحتين عشوائيتين؟
\end{exercise}

\begin{exercise}[\lvlB]
\label{ex:11:3}
\emph{ما الذي يقدر عليه منحنى واحد~\eqref{eq:nakarushton}.} (أ)~في أي مدى من السطوع تزداد الاستجابة من $10\,\%$ إلى $90\,\%$ من $r_{\max}$؟ كم رتبة هذا؟ (ب)~كم موضعًا للمنحنى المنزاح يلزم لتغطية عشر رتب من السطوع؟ واثنتي عشرة رتبة من الشدة الصوتية؟
\end{exercise}

\begin{exercise}[\lvlB]
\label{ex:11:4}
\emph{حد شفرة التوقيت في الأذن.} الفرق بين وصول الصوت إلى الأذنين يصل إلى $0.22\,\text{m}/343\,\text{m/s}$. (أ)~النبضات تسير في طور مع النغمة: حتى أي تردد يتحدد الاتجاه بها بلا لبس؟ (ب)~ترتجف النبضة بمقدار $0.5\ms$، والمطلوب تمييز $20\,\text{\textmu s}$. كم عصبونًا بمعدل $100$ نبضة في الثانية يجب أن نأخذ متوسطها خلال $50\ms$؟
\end{exercise}

\begin{exercise}[\lvlB]
\label{ex:11:5}
\emph{كاميرا أحداث على كابل العين.} كاميرا من $10^6$ بكسل بمخرج $10^7$ بت في الثانية ترسل حدثًا (العنوان والإشارة) حين يتغير $\ln I$ في البكسل بمقدار $\ln1.2$. ما أكبر سرعة يمكن أن تجري بها حافة طولها $500$ بكسل وفرق السطوع عبرها $4$ أضعاف؟ وكم إطارًا في الثانية تعطي كاميرا عادية عبر المخرج نفسه بـ$8$ بتات لكل بكسل؟
\end{exercise}

\begin{exercise}[\lvlB]
\label{ex:11:6}
\emph{نمذجة: ضبط الكسب.} يضبط المستشعر~\eqref{eq:nakarushton} القيمة $\sigma$ بـ$\tau=1\,\text{s}$ لوغاريتميًا، $\tau\,\dd\ln\sigma/\dd t=\ln I-\ln\sigma$، أو خطيًا، $\tau\,\dd\sigma/\dd t=I-\sigma$؛ والخلفية تقفز $1\to100\to1$. بعد كم ثانية تعود الاستجابة لاختبار $+10\,\%$ إلى نصف الاستجابة المتكيفة بعد القفزة صعودًا وهبوطًا، في كل من القانونين؟
\end{exercise}

\begin{exercise}[\lvlC]
\label{ex:11:7}
\emph{لأي عالم ضُبط المنحنى.} عند ضجيج مخرج ثابت صغير تكون المعلومات عن السطوع أكبر ما يمكن حين $r(I)=r_{\max}\Phi_I(I)$، حيث $\Phi_I$ دالة توزيع السطوع. (أ)~لأي توزيع يكون المنحنى~\eqref{eq:nakarushton} أمثل، وما انتشار $\log_{10}I$ فيه؟ (ب)~أي منحنى هو الأمثل عند ضجيج مخرج بواسوني؟
\end{exercise}
